\chapter{Standard Library, Examples, and Tests}

The implementation is only useful if the language can support realistic
programs. The standard library, examples, and tests are therefore part of the
artifact, not peripheral material.

\section{Standard Library}

The standard library lives under \pathcode{stdlib}. It is implemented in
\einlang{} source files and imported through the module system. Builtins are a
small fixed set known to the compiler and runtime; stdlib functions are normal
library code whenever possible.

At the time of this thesis draft, the repository contains 41 standard-library
\texttt{.ein} files. That size matters because it forces the module loader,
name resolver, type inference, monomorphization service, and backend to support
ordinary library use rather than only hand-written compiler tests.

\begin{table}[htbp]
\centering
\caption{Major standard-library areas.}
\label{tab:stdlib}
\renewcommand{\arraystretch}{1.08}
\begin{tabular}{@{}L{0.28\textwidth}L{0.60\textwidth}@{}}
\toprule
Area & Examples \\
\midrule
\texttt{std::math} & Basic arithmetic helpers, constants, trigonometric, exponential, logarithmic, hyperbolic, clamp, modulo, and special functions. \\
\texttt{std::array} & Flatten, transpose, reductions, concatenation, argmax/argmin, partition, top-k helpers, and array utility operations. \\
\texttt{std::ml} & Activations, comparison/logical ops, convolution, indexing, layers, linear algebra, normalization, pooling, reductions, recurrent ops, resizing, selection, shape ops, special ops, transforms, and utilities. \\
\texttt{std::numerics} & ODE stepping and trajectories, optimization helpers, and dynamic-programming routines. \\
\texttt{std::io} & Loading and saving NumPy arrays plus printing-oriented support. \\
\bottomrule
\end{tabular}
\end{table}

The breadth of the stdlib stresses modules, imports, function calls,
monomorphization, tensor ranks, shape operations, and backend execution.

\section{Examples}

The examples directory covers the language from small onboarding programs to
larger model-shaped workloads. Important groups include:

\begin{itemize}
\item basics and demos for syntax, arrays, functions, matrix operations, tuple
syntax, pipeline operators, and module imports;
\item autodiff examples for scalar derivatives, chains, losses, matrix
products, gradient descent, and user functions;
\item recurrence, ODE, PDE, optimization, finance, value-iteration, time-series,
and job-search examples;
\item application examples such as linear regression, decay calibration,
Kalman filtering, MNIST training, quantized MNIST inference, DeiT-Tiny, and
Whisper-Tiny.
\end{itemize}

The model examples matter because they exercise scale and shape complexity. The
MNIST examples exercise autodiff and training loops. DeiT-Tiny and Whisper-Tiny
exercise larger indexed tensor programs and standard-library ML operations.

The repository currently includes 128 \texttt{.ein} example programs. They
serve two roles: public demonstrations for users and integration pressure for
the compiler. A small syntax example usually touches parsing, name resolution,
type inference, and the NumPy expression backend. A model example additionally
touches modules, stdlib imports, shape analysis across many intermediate
tensors, lowering choices, file I/O, and backend performance diagnostics.

\section{Test Suite}

The test tree is divided by evidence type:

\begin{itemize}
\item \pathcode{tests/unit} checks narrow compiler/runtime behavior such as IR
serialization, arrays, expressions, declarations, diagnostics, type and shape
validation, pipeline validation, lazy Jacobians, recurrence storage, circular
buffers, NumPy helpers, IREE fallback behavior, and specific pass behavior.
\item \pathcode{tests/integration} checks end-to-end language behavior such as
Fibonacci recurrence, reductions, named rest patterns, recursive functions,
precision, shadowing, quantifier equivalences, scans, monomorphization, and IR
execution.
\item \pathcode{tests/stdlib} checks stdlib modules and ML operations against
NumPy references across ranks and operator families.
\item \pathcode{tests/examples} checks public examples, docs examples, MNIST
autodiff behavior, simulation accuracy, and larger demo workloads.
\end{itemize}

Reference implementations in \pathcode{tests/examples/reference_implementations.py}
provide NumPy baselines for simulation and numerical examples. Some paper
comparison sources under \pathcode{paper/compare_sources} provide nearby JAX and
PyTorch encodings for source-boundary comparison.

There are 115 Python test files in the current test tree. This count is less
important than the spread: unit tests catch small compiler invariants,
integration tests exercise end-to-end language behavior, stdlib tests compare
library functions against NumPy references, and example tests protect public
promises in the README and docs.

\section{Coverage Philosophy}

The tests are not only regression tests. They encode the claims of the system:
range and shape errors should be detected before execution; recurring bindings
should execute in legal order; bounded recurrences should use circular storage
when safe; local derivative requests should behave as shaped values; stdlib
functions should match numerical references; public examples should continue to
run as documented.
